% 프롬프트 상자 -- 프리앰블에서 % 프롬프트 상자 -- 프리앰블에서 % 프롬프트 상자 -- 프리앰블에서 % 프롬프트 상자 -- 프리앰블에서 \input{promptbox}
%
% 원문 파일을 \VerbatimInput 으로 그대로 읽는다. 논문에 옮겨 적지 않으므로 실제
% 라벨링에 쓴 문구와 갈라질 수 없다. fvextra 가 fancyvrb 를 대신한다: 프롬프트 한
% 줄이 400~580자라 줄 접기(breaklines)가 필수인데, 그 옵션은 fancyvrb 가 아니라
% fvextra 에 있다.
\usepackage[most]{tcolorbox}
\usepackage{fvextra}

\definecolor{promptrule}{HTML}{3D5A80}
\definecolor{promptback}{HTML}{F7F9FB}

% \begin{promptbox}[추가옵션]{제목} ... \end{promptbox}
\newtcolorbox{promptbox}[2][]{%
  breakable, enhanced, sharp corners,
  colback=promptback, colframe=promptrule, boxrule=0.6pt,
  left=5pt, right=5pt, top=6pt, bottom=4pt,
  fonttitle=\small\bfseries, title={#2},
  attach boxed title to top left={xshift=6pt, yshift=-2.2pt},
  boxed title style={colback=promptrule, sharp corners, boxrule=0pt,
                     left=4pt, right=4pt, top=1.5pt, bottom=1.5pt},
  #1}

% 본문: 등폭, 작게, 줄바꿈은 원문 그대로, 긴 줄은 상자 폭에서 접는다.
\fvset{fontsize=\scriptsize, fontfamily=tt, breaklines=true, breakanywhere=true,
       breakindent=0pt, xleftmargin=0pt, baselinestretch=0.95}

% 부록에서:
%   \begin{promptbox}{RoboCasa (v22)}
%     \VerbatimInput{prompts/robocasa_v22.txt}
%   \end{promptbox}

%
% 원문 파일을 \VerbatimInput 으로 그대로 읽는다. 논문에 옮겨 적지 않으므로 실제
% 라벨링에 쓴 문구와 갈라질 수 없다. fvextra 가 fancyvrb 를 대신한다: 프롬프트 한
% 줄이 400~580자라 줄 접기(breaklines)가 필수인데, 그 옵션은 fancyvrb 가 아니라
% fvextra 에 있다.
\usepackage[most]{tcolorbox}
\usepackage{fvextra}

\definecolor{promptrule}{HTML}{3D5A80}
\definecolor{promptback}{HTML}{F7F9FB}

% \begin{promptbox}[추가옵션]{제목} ... \end{promptbox}
\newtcolorbox{promptbox}[2][]{%
  breakable, enhanced, sharp corners,
  colback=promptback, colframe=promptrule, boxrule=0.6pt,
  left=5pt, right=5pt, top=6pt, bottom=4pt,
  fonttitle=\small\bfseries, title={#2},
  attach boxed title to top left={xshift=6pt, yshift=-2.2pt},
  boxed title style={colback=promptrule, sharp corners, boxrule=0pt,
                     left=4pt, right=4pt, top=1.5pt, bottom=1.5pt},
  #1}

% 본문: 등폭, 작게, 줄바꿈은 원문 그대로, 긴 줄은 상자 폭에서 접는다.
\fvset{fontsize=\scriptsize, fontfamily=tt, breaklines=true, breakanywhere=true,
       breakindent=0pt, xleftmargin=0pt, baselinestretch=0.95}

% 부록에서:
%   \begin{promptbox}{RoboCasa (v22)}
%     \VerbatimInput{prompts/robocasa_v22.txt}
%   \end{promptbox}

%
% 원문 파일을 \VerbatimInput 으로 그대로 읽는다. 논문에 옮겨 적지 않으므로 실제
% 라벨링에 쓴 문구와 갈라질 수 없다. fvextra 가 fancyvrb 를 대신한다: 프롬프트 한
% 줄이 400~580자라 줄 접기(breaklines)가 필수인데, 그 옵션은 fancyvrb 가 아니라
% fvextra 에 있다.
\usepackage[most]{tcolorbox}
\usepackage{fvextra}

\definecolor{promptrule}{HTML}{3D5A80}
\definecolor{promptback}{HTML}{F7F9FB}

% \begin{promptbox}[추가옵션]{제목} ... \end{promptbox}
\newtcolorbox{promptbox}[2][]{%
  breakable, enhanced, sharp corners,
  colback=promptback, colframe=promptrule, boxrule=0.6pt,
  left=5pt, right=5pt, top=6pt, bottom=4pt,
  fonttitle=\small\bfseries, title={#2},
  attach boxed title to top left={xshift=6pt, yshift=-2.2pt},
  boxed title style={colback=promptrule, sharp corners, boxrule=0pt,
                     left=4pt, right=4pt, top=1.5pt, bottom=1.5pt},
  #1}

% 본문: 등폭, 작게, 줄바꿈은 원문 그대로, 긴 줄은 상자 폭에서 접는다.
\fvset{fontsize=\scriptsize, fontfamily=tt, breaklines=true, breakanywhere=true,
       breakindent=0pt, xleftmargin=0pt, baselinestretch=0.95}

% 부록에서:
%   \begin{promptbox}{RoboCasa (v22)}
%     \VerbatimInput{prompts/robocasa_v22.txt}
%   \end{promptbox}

%
% 원문 파일을 \VerbatimInput 으로 그대로 읽는다. 논문에 옮겨 적지 않으므로 실제
% 라벨링에 쓴 문구와 갈라질 수 없다. fvextra 가 fancyvrb 를 대신한다: 프롬프트 한
% 줄이 400~580자라 줄 접기(breaklines)가 필수인데, 그 옵션은 fancyvrb 가 아니라
% fvextra 에 있다.
\usepackage[most]{tcolorbox}
\usepackage{fvextra}

\definecolor{promptrule}{HTML}{3D5A80}
\definecolor{promptback}{HTML}{F7F9FB}

% \begin{promptbox}[추가옵션]{제목} ... \end{promptbox}
\newtcolorbox{promptbox}[2][]{%
  breakable, enhanced, sharp corners,
  colback=promptback, colframe=promptrule, boxrule=0.6pt,
  left=5pt, right=5pt, top=6pt, bottom=4pt,
  fonttitle=\small\bfseries, title={#2},
  attach boxed title to top left={xshift=6pt, yshift=-2.2pt},
  boxed title style={colback=promptrule, sharp corners, boxrule=0pt,
                     left=4pt, right=4pt, top=1.5pt, bottom=1.5pt},
  #1}

% 본문: 등폭, 작게, 줄바꿈은 원문 그대로, 긴 줄은 상자 폭에서 접는다.
\fvset{fontsize=\scriptsize, fontfamily=tt, breaklines=true, breakanywhere=true,
       breakindent=0pt, xleftmargin=0pt, baselinestretch=0.95}

% 부록에서:
%   \begin{promptbox}{RoboCasa (v22)}
%     \VerbatimInput{prompts/robocasa_v22.txt}
%   \end{promptbox}
